\chapter{The Forces from the Field}
\label{chap:forces}

If matter is torsitons, the forces between them are not separate ingredients either: they are
what the same gravity--torsion field does when two knots overlap and exchange. This chapter
walks the strong and electroweak interactions back to properties of that one field, and asks
the hard question of the programme---whether the electroweak precision parameter $S$ can be
brought below its measured bound. The answer decides whether Part~III lives or retreats to
Part~II.

\section{Colour is forced by the label}
\label{sec:colour}

A torsiton built from a Dirac field carries the Pauli exclusion label, and for three filled
states that label is three-valued. Treating it as an internal $SU(3)$ charge is not a choice
bolted on---it is forced by the requirement that the bound state antisymmetrise. The resulting
channel factors $\langle T_1\!\cdot\!T_2\rangle$ are then pure group theory, and they
\emph{are} those of QCD: the singlet attractive ($-4/3$), the octet repulsive ($+1/6$), the
antitriplet ($-2/3$), the sextet ($+1/3$). Mesons (singlet) and baryons (three colours) bind;
free quarks and free diquarks do not. The colour structure of the strong interaction is
recovered, not assumed (Figure~\ref{fig:color}).

\begin{figure}[t]
  \centering
  \includegraphics[width=\textwidth]{color_force}
  \caption{Colour channel factors from the $SU(3)$ Casimirs forced by the Pauli label. The
  attractive singlet and the repulsive octet reproduce QCD's channel structure; every
  perturbative overlap channel screens to zero, so confinement must come from elsewhere
  (Section~\ref{sec:confinement}).}
  \label{fig:color}
\end{figure}

\section{The four-fermion term, Fierzed}
\label{sec:fierz}

The torsion-induced interaction is a single axial--axial four-fermion term \citep{hehl1976}. A
Fierz transformation re-expresses it in the scalar, pseudoscalar, vector, axial, and tensor
channels; carried out explicitly with the Dirac matrices it gives equal vector and axial
couplings, $G_A = G_V$ (Figure~\ref{fig:fierz}). This is not a cosmetic detail: a sector with
$G_A=G_V$ is precisely the ``walking-favourable'' case, the one whose vector and axial current
correlators tend to equalise---which, as Section~\ref{sec:S} shows, is exactly what is needed
to push the electroweak $S$ down.

\begin{figure}[t]
  \centering
  \includegraphics[width=\textwidth]{fierz}
  \caption{The Hehl--Datta four-fermion term, Fierzed into channels: the vector and axial
  couplings come out equal, $G_A=G_V$---the walking-favourable starting point for the
  electroweak parameter $S$.}
  \label{fig:fierz}
\end{figure}

\section{Confinement: one substrate, two regimes}
\label{sec:confinement}

The overlap (one-gluon-exchange) forces of Section~\ref{sec:colour} all \emph{screen}: they
fall to zero at large separation. Confinement---a potential that rises forever, $V(L)=\sigma
L$---is not produced by any pairwise overlap; it is a non-perturbative property of the
non-abelian field. The candidate mechanism is the dual superconductor: if the chiral condensate
expels colour-electric flux, it squeezes that flux into a Nielsen--Olesen tube of finite width,
whose energy grows linearly with length (Figure~\ref{fig:flux}). At the Bogomolny point the
string tension saturates the topological bound, $\sigma = 2\pi v^2$---so the \emph{same}
condensate that sets configurational mass sets the confinement scale. One substrate does both
jobs; the potential is the Cornell sum $V(r)\approx -\kappa/r + \sigma L$, the overlap dressing
the short distance, the substrate confining the long. The Bogomolny value is not assumed: solving
the winding-1 Abrikosov--Nielsen--Olesen vortex of the condensate gives $\sigma = 2\pi v^2\,
\varepsilon(\beta)$ with $\beta=(m_{\mathrm{scalar}}/m_{\mathrm{vector}})^2$, and the vortex BVP
returns $\varepsilon(1)=1$ exactly at the BPS point (type-I $\varepsilon<1$, type-II $\varepsilon>1$;
\texttt{sims/.../l2\_string\_tension.py}). So $\sigma=2\pi v^2$ holds precisely \emph{when the
vacuum is BPS}, and the residual factor is fixed by $\beta$ alone. The area law itself is already
in hand---a measured $\sigma>0$ from Wilson loops on pure $SU(3)$. What remains is whether the
gravity--torsion vacuum actually realises dual superconductivity (the dual order parameter), the
cleanest lattice question (Appendix), and $\beta$, which sets whether the saturation is exact.

\begin{figure}[t]
  \centering
  \includegraphics[width=\textwidth]{flux_tube}
  \caption{A confining flux tube from the condensate. A Nielsen--Olesen vortex squeezes
  colour-electric flux into a string of tension $\sigma=2\pi v^2$ (the same condensate scale
  that sets mass), giving the linear confinement the screening overlaps cannot. Conditional on
  the vacuum being a dual superconductor.}
  \label{fig:flux}
\end{figure}

\section{The electroweak $S$: the make-or-break}
\label{sec:S}

If electroweak symmetry is broken by the gravity--torsion condensate rather than an elementary
Higgs, the theory is technicolor-like \citep{weinberg1979,susskind1979}, and its central danger
is the Peskin--Takeuchi parameter $S$ \citep{peskin1992}. A QCD-like sector gives $S\approx0.25$,
well above the LEP bound $S\lesssim0.1$: the ``technicolor graveyard'' (Figure~\ref{fig:electroweak-S}).
Two facts move the theory
off that headstone. First, the $G_A=G_V$ of Section~\ref{sec:fierz} equalises the vector and
axial correlators, and a random-phase (RPA) resummation of the simpler numerical model shows the
ratio $S_{\mathrm{torsion}}/S_{\mathrm{QCD}}$ falling as the sector walks
(Figure~\ref{fig:rpa})---a real, growing relative advantage. Second, near-conformal (``walking'')
dynamics \citep{holdom1981,dietrich2007} raises the resonance masses relative to $f_\pi$, and the
Weinberg sum rules \citep{weinberg1967} give
\[
  S = 4\pi\left(\frac{f_\pi}{M_V}\right)^2\!\left(1+\frac{M_V^2}{M_A^2}\right),
\]
which falls below $0.1$ once $M_V/f_\pi \gtrsim 13$ (versus QCD's $\sim 8$). The simpler numerical
work---analytic Fierz algebra, RPA moments, sum-rule saturation, as opposed to a full lattice
treatment---places the theory \emph{favourably}, but does not settle the absolute value: whether
the walking sector actually clears $S<0.1$ is a frontier lattice computation
(Appendix~\ref{app:lattice}), and the field has not settled the analogous technicolor question in
forty years.
We bank on it provisionally and mark it as the load-bearing assumption.

\begin{figure}[t]
  \centering
  \includegraphics[width=\textwidth]{electroweak_S}
  \caption{The electroweak $S$. A QCD-like sector sits in the graveyard ($S\sim0.25$); the
  $G_A=G_V$ torsion sector, RPA-resummed, gains a growing relative advantage as it walks. The
  direction is favourable; the absolute $S<0.1$ is the lattice make-or-break.}
  \label{fig:electroweak-S}
\end{figure}

\begin{figure}[t]
  \centering
  \includegraphics[width=\textwidth]{rpa}
  \caption{RPA resummation: with $G_A=G_V$ the torsion sector's $S$ falls relative to a
  QCD-like sector as the dynamics walk---a genuine, growing advantage (shown as a ratio; the
  absolute value is owed to the lattice).}
  \label{fig:rpa}
\end{figure}

\section{$W$, $Z$, and the composite Higgs}
\label{sec:wzh}

Banking on $S$, the boson masses are then the cleanest numbers in the Part---though it is worth
being precise about what is \emph{computed} and what is \emph{input}. We take the condensate scale
$v=246~\mathrm{GeV}$ (in the framework this is the infrared scale $\Lambda$ of
Section~\ref{sec:ir-potential}, generated rather than an inserted Higgs vacuum value, but its
\emph{numerical} value is fixed to experiment here) and the measured gauge couplings $g,g'$, so
that
\[
  m_W = \tfrac12 g v, \qquad m_Z = \tfrac12\sqrt{g^2+g'^2}\,v
\]
are not predictions of $80$ and $91~\mathrm{GeV}$ from nothing---they reproduce the Standard
Model's own tree-level relations, with the same three inputs. These are therefore not ab-initio
numbers. What \emph{is} a genuine, input-free prediction is the structure: a custodial (doublet)
condensate gives $\rho = m_W^2/(m_Z^2\cos^2\theta_W)=1$, equivalently $m_W=m_Z\cos\theta_W$---one
fewer parameter than a general Higgs sector, observed to $\sim0.1\%$ (Figure~\ref{fig:ewmasses}).
The $125~
\mathrm{GeV}$ scalar is then \emph{composite}: a torsiton--anti-torsiton bound state, not an
elementary field. Its mass is a strong-dynamics bound state---owed to the lattice---but its
\emph{lightness} is explained: a generic composite scalar would be heavy ($\sim$~TeV, the
excluded heavy-sigma), whereas a walking sector yields a light pseudo-dilaton at $\sim0.5v$,
exactly where the observed Higgs sits. The Higgs is light because the sector walks---the same
walking the $S$ parameter needs.

This is worth stating bluntly, because it is the framework's most direct collision with
established particle physics. The 2013 Nobel Prize honoured the Higgs \emph{mechanism}---
electroweak symmetry broken by a condensate---and that mechanism this framework keeps entirely;
the bounce-condensate breaks the symmetry exactly as advertised. What it does \emph{not} keep is
the elementary Higgs \emph{field}. In this reading the $125~\mathrm{GeV}$ particle is not a
fundamental scalar but the lightest composite of the new strong sector---the analogue of the
$\sigma$/$f_0(500)$ meson of QCD, a torsiton--anti-torsiton state. The discovery at the LHC
confirmed the \emph{mechanism}; it did not establish that the boson is elementary, and the
experiments that could tell an elementary scalar from a composite one---precision coupling
measurements, off-shell behaviour, the search for partners---are not yet decisive. To that
extent the prize was awarded on the mechanism while the deeper question, what the boson
\emph{is}, remained open.

And the framework makes that question answerable, because a composite Higgs cannot be alone. If
the $125~\mathrm{GeV}$ scalar is a torsiton pair bound by the new strong force, then---exactly as
QCD has not one meson but a whole spectrum ($\sigma$, $\rho$, $a_1$, $\dots$)---there must be a
\emph{tower} of further composites: heavier scalars, pseudoscalars, and especially vector
resonances (the ``techni-$\rho$'') built from torsiton pairs, plus the technipions already eaten
to make the $W$ and $Z$ massive. These are short-lived torsiton--anti-torsiton (and multi-torsiton)
states at the TeV scale, and finding even one would settle the matter. So the sharp,
falsifiable prediction that distinguishes this picture from the Standard Model is twofold: the
$125~\mathrm{GeV}$ boson should eventually betray its compositeness (form factor, coupling
deviations) where the Standard Model says ``point-like forever''; and it should not be the only
Higgs-like state---a spectrum of heavier composite resonances should appear as the energy reach
climbs. The Standard Model predicts a desert above the Higgs; this framework predicts company.

\begin{figure}[t]
  \centering
  \includegraphics[width=\textwidth]{electroweak_masses}
  \caption{$W$, $Z$, and Higgs in the composite/walking reading. $m_W,m_Z$ follow from the
  condensate scale and the gauge couplings (within $0.4\%$), with custodial $\rho=1$ a genuine
  prediction. The Higgs is a composite scalar; the observed $125~\mathrm{GeV}$ sits at the
  light, walking-dilaton edge, far from the excluded heavy-sigma.}
  \label{fig:ewmasses}
\end{figure}

There is no separate generation-changing boson to find: generations are internal rungs of one
torsiton, so a generation transition is an internal de-excitation carried by the \emph{existing}
weak current, and the CKM/PMNS mixings are the wavefunction overlaps between rungs. The place
this picture is tested is the flavour-changing neutral currents that plague extended
technicolor---the experimental front where the parsimony lives or dies.
